\chapter{Mesure de la performance}\label{ch:rs:performance-measurement}

Deux stratégies affichent des ratios de Sharpe de 1.00 et 1.09 sur cinq ans de rendements quotidiens. La première a perdu 33.6~\% de son sommet à son creux et est restée 791 jours de
bourse, soit plus de trois ans, sous l'eau. La seconde vend des options de vente hors de la monnaie et n'a jamais connu de trimestre perdant~; son pire drawdown était de 14.8~\%.
Les deux historiques ont été tirés des simulations du chapitre, portant sur des stratégies de volatilités de long terme de 15 et 17~\%, et la seconde a été choisie
parmi les 9~\% de ses historiques simulés sur cinq ans sans trimestre perdant. Sur les cinq années suivantes, la première a une probabilité de 14~\% d'avoir un drawdown de
30~\% ou plus, la seconde, 41~\%. Le ratio de Sharpe est le bon premier chiffre et un dangereux dernier chiffre. Ce chapitre construit le reste de la fiche
de synthèse (erreurs types, drawdowns, \omterm{def:rs:vectorised-backtests:turnover}{rotation du portefeuille}, forme des résultats possibles, indice de référence) sous forme de \texttt{firm.perf}, qui lit le \texttt{BacktestResult} des
chapitres 16 à 18, et le lit pour quatre flux synthétiques et le portefeuille momentum du chapitre~16.

\section{Le ratio de Sharpe et son erreur type}

Le ratio de Sharpe (livre~4, chapitre~11) est le rendement excédentaire moyen rapporté à son écart type~; Sharpe l'a introduit en 1966 comme ratio reward-to-variability
et lui a donné son nom actuel en 1994.

\begin{definition}[Annualisation]\label{def:rs:performance-measurement:annualisation}
L'\emph{annualisation}\index{annualisation} convertit une statistique mesurée sur des périodes d'une durée donnée en valeur annuelle~: avec $q$ périodes par an, la moyenne est multipliée par $q$,
la volatilité par $\sqrt q$ et le ratio de Sharpe par $\sqrt q$. La mise à l'échelle de la volatilité et du ratio de Sharpe n'est valable que si les rendements sont
sans autocorrélation sérielle.
\end{definition}

L'erreur type d'un ratio de Sharpe estimé, pour des rendements normaux indépendants, vaut environ $\sqrt{(1 + \mathrm{SR}^2/2)/T}$ par période sur $T$ périodes (Lo,
2002~; livre~4, chapitre~11). Annualisée, elle vaut presque $\sqrt{q/T}$~: la durée de l'historique en années la fixe, et non la fréquence. Cinq ans de rendements quotidiens
donnent une erreur type de 0.45, quelle que soit la stratégie~; un ratio de Sharpe mesuré de 1.0 a un intervalle à 95~\% allant de 0.12 à 1.88. Les quatre flux de cinq ans du chapitre
ont cette erreur type iid (la fiche de synthèse du tutoriel)~; l'erreur type HAC, qui tient aussi compte de l'autocorrélation sérielle et des queues épaisses, est proche pour les
flux de tendance et de volatilité courte, plus grande pour les journées à queues épaisses du teneur de marché (0.52), et maximale pour un portefeuille lissé.

\begin{definition}[Lissage des rendements, ratio de Sharpe corrigé de l'autocorrélation]\label{def:rs:performance-measurement:smoothing}
Le \emph{lissage des rendements}\index{lissage des rendements} est la publication d'une moyenne mobile des rendements vrais, $r^o_t = \sum_{j=0}^{k}\theta_j r_{t-j}$ avec
$\theta_j \ge 0$ et $\sum_j\theta_j = 1$, comme c'est le cas lorsque des actifs illiquides sont valorisés à des cours périmés ou expertisés (Getmansky, Lo et Makarov). Le
\emph{ratio de Sharpe corrigé de l'autocorrélation}\index{ratio de Sharpe corrigé de l'autocorrélation} annualise un ratio de Sharpe par période $\mathrm{SR}_1$ à l'aide des
autocorrélations $\rho_k$ des rendements~: $\mathrm{SR}(q) = \mathrm{SR}_1\, q\big/\sqrt{q + 2\sum_{k=1}^{q-1}(q-k)\rho_k}$ (Lo, 2002).
\end{definition}

\begin{proposition}[Effet du lissage]\label{prop:rs:performance-measurement:smoothing}
Si les rendements vrais sont indépendants, de moyenne $\mu$ et de variance $\sigma^2$, les rendements publiés, lissés avec le profil $\theta$, ont une moyenne $\mu$, une variance
$\sigma^2\sum_j\theta_j^2$ et des autocorrélations $\rho_k = \sum_j\theta_j\theta_{j+k}\big/\sum_j\theta_j^2$. Le ratio de Sharpe naïf est gonflé d'un facteur
$1/\sqrt{\sum_j\theta_j^2}$~; le ratio corrigé de l'autocorrélation, calculé avec les autocorrélations théoriques, supprime ce gonflement à des termes d'ordre relatif
$1/q$ près.
\end{proposition}

\begin{proof}
La moyenne et les autocovariances $\sigma^2\sum_j\theta_j\theta_{j+k}$ découlent de la linéarité et de l'indépendance. Sur $q$ périodes, les rendements publiés ont la même somme que les rendements vrais
sauf aux bords, si bien que leur variance vaut $q\sigma^2$ à des termes d'ordre un près~; dans la formule de Lo, le dénominateur $q + 2\sum(q-k)\rho_k$ multiplié par
$\sigma^2\sum\theta_j^2$ est exactement la variance de la somme de $q$ rendements publiés, soit $q\sigma^2$ augmentée des termes de bord.
\end{proof}

Avec le profil $(0.5, 0.3, 0.2)$, la variance est multipliée par 0.38 et la volatilité par 0.616~; un ratio de Sharpe vrai de 0.6 est publié comme 0.97, avec
des premières autocorrélations de 0.55 et 0.26. Le portefeuille illiquide du chapitre a exactement ce profil, et son historique sur cinq ans en montre l'effet~:

\begin{center}
\small
\begin{tabular}{@{}lcc@{}}
\toprule
le portefeuille lissé, 60 mois & publié & vrai ou corrigé\\
\midrule
volatilité annuelle & 7.4~\% & 11.2~\% (vraie), 12.7~\% (non lissée)\\
ratio de Sharpe & 0.99 & 0.66 (vrai), 0.61 (Lo), 0.59 (non lissé)\\
erreur type du ratio de Sharpe & 0.46 (iid) & 0.63 (HAC)\\
deux premières autocorrélations & 0.64, 0.25 & \\
profil de lissage & & $(0.43, 0.36, 0.21)$ estimé, $(0.5, 0.3, 0.2)$ vrai\\
\bottomrule
\end{tabular}
\end{center}

Trois corrections concordent~: l'\omterm{def:rs:performance-measurement:annualisation}{annualisation} de Lo avec les autocorrélations empiriques, la série non lissée (la moyenne mobile inversée avec le profil estimé
à partir des deux premières autocorrélations, voir le \cref{lst:rs:performance-measurement:unsmooth}), et l'erreur type HAC, qui indique que le 0.99 publié est aussi
moins précis qu'il n'y paraît. Lo a montré que les ratios de Sharpe des hedge funds étaient surestimés jusqu'à 65~\% à cause de l'autocorrélation sérielle~; Getmansky, Lo et
Makarov l'ont imputé à l'illiquidité.

\section{Drawdowns}

\begin{definition}[Drawdown maximal, durée du drawdown, ratio de Calmar]\label{def:rs:performance-measurement:drawdown}
Le \emph{drawdown maximal}\index{drawdown maximal} d'un historique est sa plus forte chute de richesse depuis un sommet courant, exprimée en fraction de ce sommet (drawdown~: livre~2,
chapitre~29). La \emph{durée du drawdown}\index{durée du drawdown} d'un épisode est le temps écoulé entre le sommet et le recouvrement de ce sommet, ou jusqu'à la fin de
l'historique s'il n'a pas recouvré ce sommet. Le \emph{ratio de Calmar}\index{ratio de Calmar} est le rendement annualisé rapporté au drawdown maximal en valeur absolue, traditionnellement calculé sur trois ans.
\end{definition}

Les drawdowns sont ce que les investisseurs subissent et ce qui met fin aux stratégies~; ce sont aussi les chiffres les moins stables de la fiche de synthèse, car un seul chemin donne un seul
maximum. Pour un mouvement brownien à dérive positive, le \omterm{def:rs:performance-measurement:drawdown}{drawdown maximal} espéré ne croît que logarithmiquement avec l'horizon (Magdon-Ismail, Atiya, Pratap et
Abu-Mostafa)~; le maximum d'un historique isolé est un seul tirage dans une large distribution. Le 33.6~\% sur cinq ans du flux de tendance est un mauvais tirage~: son modèle lui
donne une probabilité de 14~\% de subir un drawdown de 30~\% ou plus en cinq ans. Le 14.8~\% du flux de volatilité courte est un bon tirage~: sa probabilité est de 41~\%. Sur 200 années simulées, les deux atteignent des
drawdowns maximaux de 37~\% et 72~\%, avec des ratios de Sharpe de 0.83 et 0.62 (\cref{fig:rs:performance-measurement:paths,fig:rs:performance-measurement:ddprob}). Les \omterm{def:rs:performance-measurement:drawdown}{ratios de Calmar}
des historiques affichés, 0.43 et 0.78, les classent à l'envers.

\begin{omfigure}
\begin{tikzpicture}
\begin{axis}[name=top,width=12cm,height=5.2cm,ylabel={\small richesse},tick label style={font=\footnotesize},xmin=0,xmax=5,ymin=0.85,ymax=2.9,
  grid=major,grid style={gray!25},xticklabels={},legend columns=4,
  legend style={font=\footnotesize,at={(0.5,1.03)},anchor=south,draw=none,fill=none,/tikz/every even column/.append style={column sep=6pt}},
  table/col sep=comma]
\addplot[black!60,thick] table[x=year,y=trend]{figdata/research/22-performance-measurement/paths.csv};
\addplot[omThm,thick] table[x=year,y=shortvol]{figdata/research/22-performance-measurement/paths.csv};
\addplot[omDef,thick] table[x=year,y=smoothed]{figdata/research/22-performance-measurement/paths.csv};
\addplot[green!45!black,thick,dashed] table[x=year,y=mm]{figdata/research/22-performance-measurement/paths.csv};
\legend{tendance,volatilité courte,portefeuille lissé,teneur de marché}
\end{axis}
\begin{axis}[at={(top.south west)},anchor=north west,yshift=-0.2cm,width=12cm,height=4.2cm,xlabel={\small années},ylabel={\small drawdown},
  tick label style={font=\footnotesize},xmin=0,xmax=5,ymin=-0.36,ymax=0.02,grid=major,grid style={gray!25},
  yticklabel style={/pgf/number format/fixed,/pgf/number format/precision=2},scaled y ticks=false,table/col sep=comma]
\addplot[black!60,thick] table[x=year,y=dd_trend]{figdata/research/22-performance-measurement/paths.csv};
\addplot[omThm,thick] table[x=year,y=dd_shortvol]{figdata/research/22-performance-measurement/paths.csv};
\addplot[omDef,thick] table[x=year,y=dd_smoothed]{figdata/research/22-performance-measurement/paths.csv};
\addplot[green!45!black,thick,dashed] table[x=year,y=dd_mm]{figdata/research/22-performance-measurement/paths.csv};
\end{axis}
\end{tikzpicture}

\omcaption{Les quatre historiques de cinq ans affichés~: richesse (en haut) et drawdown depuis le sommet courant (en bas). Le portefeuille lissé est mensuel et tracé en paliers
d'un mois. Données~: \texttt{rs\_perf.histories}.}\label{fig:rs:performance-measurement:paths}
\end{omfigure}

\begin{omfigure}
\begin{tikzpicture}
\begin{axis}[width=10.5cm,height=5.4cm,xlabel={horizon (années)},ylabel={P(drawdown $\ge$ 30~\%)},tick label style={font=\small},label style={font=\small},
  xmin=0.5,xmax=10.5,ymin=0,ymax=0.7,ybar,bar width=8pt,grid=major,grid style={gray!25},xtick={1,2,3,4,5,6,7,8,9,10},legend columns=2,
  yticklabel style={/pgf/number format/fixed,/pgf/number format/precision=1},scaled y ticks=false,
  legend style={font=\footnotesize,at={(0.5,-0.3)},anchor=north,draw=none,fill=none,/tikz/every even column/.append style={column sep=8pt}},
  table/col sep=comma]
\addplot[fill=gray!45,draw=gray] table[x=years,y=trend]{figdata/research/22-performance-measurement/ddprob.csv};
\addplot[fill=omThm!55,draw=omThm] table[x=years,y=shortvol]{figdata/research/22-performance-measurement/ddprob.csv};
\legend{tendance,volatilité courte}
\end{axis}
\end{tikzpicture}

\omcaption{La probabilité d'un drawdown de 30~\% ou plus à l'intérieur d'un horizon, sur 4\,000 trajectoires simulées du modèle de chaque stratégie (volatilités de long terme de 15~\% et
17~\%). Données~: \texttt{rs\_perf.dd\_probability}.}\label{fig:rs:performance-measurement:ddprob}
\end{omfigure}

\section{Rotation, taux de réussite et durée de détention}

\begin{definition}[Taux de réussite, profit factor]\label{def:rs:performance-measurement:hit}
Le \emph{taux de réussite}\index{taux de réussite} (hit rate) d'un flux de rendements est la part de ses périodes non nulles dont le rendement est positif. Son \emph{profit factor}\index{profit factor}
est la somme de ses gains divisée par la somme de ses pertes.
\end{definition}

Un \omterm{def:rs:performance-measurement:hit}{taux de réussite} indique la fréquence, non l'ampleur~: le flux de volatilité courte gagne 76~\% des jours et le flux de tendance 53~\%, mais le second a la meilleure queue de distribution. Un
\omterm{def:rs:performance-measurement:hit}{profit factor} supérieur à un signale seulement une espérance positive et rien de plus~; il égale le \omterm{def:rs:performance-measurement:sortino}{ratio Oméga} en zéro (section suivante), puisque les deux sont le gain moyen divisé par la perte moyenne.
Les 65~\% de jours gagnants et le \omterm{def:rs:performance-measurement:hit}{profit factor} de 2.07 du teneur de marché sont ceux d'une activité qui capte une fourchette avec de rares jours défavorables~; son ratio de Sharpe de 4.04
(erreur type 0.45, HAC 0.52) est réel et ne dit rien de la capacité (chapitre~28).

La rotation (chapitre~16) fixe les coûts et, avec l'exposition brute, fixe la durée moyenne de détention~: un livre d'exposition brute $G$ qui échange $T$ du capital dans un
seul sens par période se renouvelle entièrement tous les $G/T$ périodes. Le livre momentum du chapitre~16, reconstruit chaque mois sur les 500 noms les plus liquides, échange 1.86~\% du capital par jour et
détient une position pendant 48 jours de bourse en moyenne~; son ratio de Sharpe après coûts vaut 0.17 avec une erreur type de 0.32 sur dix ans, et son drawdown le plus long a duré 2\,079 de ses
2\,520 jours. La durée de détention d'une stratégie doit correspondre à la demi-vie de son signal (chapitre~13)~; une durée de détention nettement inférieure à la demi-vie
est une rotation payée pour du bruit.

\section{La forme des issues}

\begin{definition}[Ratio de Sortino, ratio Oméga]\label{def:rs:performance-measurement:sortino}
Le \emph{ratio de Sortino}\index{ratio de Sortino} (Sortino ratio) est le rendement moyen en excès d'un rendement minimal acceptable, divisé par la déviation à la baisse, c'est-à-dire la racine quadratique moyenne
des écarts sous ce seuil (Sortino et Price). Le \emph{ratio Oméga}\index{ratio Oméga} (Omega ratio) au seuil $L$ vaut $\mathbb{E}[(r-L)^+]\,/\,\mathbb{E}[(L-r)^+]$,
soit le gain espéré au-dessus du seuil divisé par la perte espérée en dessous (Keating et Shadwick)~; il vaut un à la moyenne.
\end{definition}

Les mesures de risque baissier ont été conçues pour cesser de pénaliser les stratégies pour leur côté haussier. Elles ne voient pas une baisse que l'échantillon ne contient pas. Dans l'historique affiché
les rendements quotidiens du flux de volatilité courte ont une asymétrie de $+5.38$ et un aplatissement de 102.5~; sur 100 années simulées, l'asymétrie vaut $-1.23$ et l'aplatissement
121. Les cinq années retenues ne comportent que des rebonds, pas de krachs, et son \omterm{def:rs:performance-measurement:sortino}{ratio de Sortino} (1.87) et son \omterm{def:rs:performance-measurement:sortino}{ratio Oméga} (1.49) dépassent ceux du flux de tendance (1.48 et 1.17). Vendre des options
fait monter le ratio de Sharpe et les ratios de risque baissier d'un historique qui n'a pas encore rencontré le paiement des options~; Goetzmann, Ingersoll, Spiegel et Welch ont montré comment
chacune des mesures usuelles peut être manipulée de la sorte, et ont caractérisé une mesure qui ne peut pas l'être~: une moyenne d'une utilité en puissance calculée sur l'historique des rendements.

La parade consiste à mesurer la stratégie et non l'historique~: savoir ce qui lui manque, la simuler ou la soumettre à des tests de résistance là où l'historique se tait (ceux du
livre~6), et lire les statistiques de forme comme des descriptions d'un échantillon.

\section{Mesures relatives à un indice de référence}

\begin{definition}[Bêta de marché]\label{def:rs:performance-measurement:beta}
Le \emph{bêta de marché}\index{bêta de marché} (market beta) d'une stratégie est la pente de ses rendements régressés sur ceux du marché, $\beta = \operatorname{Cov}(r, r_m)/\operatorname{Var}(r_m)$,
estimé par régression avec une ordonnée à l'origine, c'est-à-dire l'alpha (livre~1, chapitre~1)~; les erreurs types sont HAC.
\end{definition}

Régressé sur l'indice sur lequel il vend des options de vente, le flux de volatilité courte de l'historique affiché a un bêta de 0.45 (erreur type 0.09) et un alpha annuel de 5.6~\%
($t = 1.38$). Un bêta est un résumé linéaire d'une exposition non linéaire. Estimés séparément sur les jours de baisse et les jours de hausse de l'indice, les bêtas de l'historique valent 0.59 et
0.54~; sur 100 années simulées, ils valent 0.98 et 0.71, et les rendements mensuels reproduisent le paiement d'une option de vente vendue (la \cref{fig:rs:performance-measurement:scatter})~:
plat au-dessus, pentu en dessous. Le bêta du flux de tendance vaut 0.00 (0.03), avec une statistique $t$ d'alpha de 2.15 sur cinq ans~; celui du teneur de marché vaut 0.04 (0.01), $t =
9.1$. Face à son indice de référence, un gérant rapporte le ratio d'information et le tracking error (livre~1, chapitre~3)~; la même prudence s'impose.

\begin{omfigure}
\begin{tikzpicture}
\begin{axis}[width=9cm,height=6cm,xlabel={rendement mensuel de l'indice},ylabel={rendement mensuel de la volatilité courte},tick label style={font=\small},label style={font=\small},
  xmin=-0.19,xmax=0.14,ymin=-0.45,ymax=0.2,grid=major,grid style={gray!25},
  xticklabel style={/pgf/number format/fixed,/pgf/number format/precision=2},yticklabel style={/pgf/number format/fixed,/pgf/number format/precision=1},
  scaled ticks=false,table/col sep=comma]
\addplot[only marks,mark=*,mark size=0.7pt,omThm] table[x=index,y=shortvol]{figdata/research/22-performance-measurement/scatter.csv};
\end{axis}
\end{tikzpicture}

\omcaption{Rendements mensuels de la stratégie de volatilité courte face à l'indice, sur 100 années simulées~: le paiement d'une option de vente vendue, avec des bêtas quotidiens de 0.98 les
jours de baisse et de 0.71 les jours de hausse. Données~: \texttt{rs\_perf.monthly\_scatter}.}\label{fig:rs:performance-measurement:scatter}
\end{omfigure}

\section{Tutoriel~: quatre flux et un livre}\label{tut:rs:performance-measurement:tut}

\begin{tutorial}
\textbf{Objectif.} Calculer la fiche de synthèse de \texttt{firm.perf} pour les quatre flux synthétiques et le livre momentum du chapitre~16, et lire ce que chaque chiffre dissimule.
\textbf{État final~:} le tableau de synthèse ci-dessous, les \cref{fig:rs:performance-measurement:paths,fig:rs:performance-measurement:ddprob,fig:rs:performance-measurement:scatter}.
\begin{tutsteps}
\item \textbf{\omterm{def:rs:performance-measurement:annualisation}{Annualisation} de Lo}~: le ratio de Sharpe par période multiplié par $q/\sqrt{q + 2\sum(q-k)\rho_k}$.
\omcode{code/firm/perf/firm_perf.py}{63}{72}{Le ratio de Sharpe corrigé de l'autocorrélation.}\label{lst:rs:performance-measurement:lo}
\item \textbf{Drawdowns et leurs épisodes}.
\omcode{code/firm/perf/firm_perf.py}{79}{105}{Drawdown, drawdown maximal et épisodes de drawdown.}\label{lst:rs:performance-measurement:dd}
\item \textbf{Délissage}~: inverser la moyenne mobile avec le profil estimé.
\omcode{code/firm/perf/firm_perf.py}{175}{183}{Reconstitution des vrais rendements à partir des rendements lissés.}\label{lst:rs:performance-measurement:unsmooth}
\item \textbf{Exécuter} \texttt{rs\_perf.histories()}, \texttt{stream\_results()}, \texttt{momentum()} et \texttt{tear\_sheet} sur chacun, puis \texttt{smoothed\_report()},
\texttt{dd\_probability()} et \texttt{fig\_perf.py}. Les fiches de synthèse portent sur cinq ans de rendements quotidiens (le livre lissé~: 60 mois) et sur les
dix ans du livre momentum, avec les bêtas par rapport à l'indice de la simulation de volatilité courte et, pour le momentum, au marché synthétique~:
\end{tutsteps}
\begin{center}
\small
\setlength{\tabcolsep}{4pt}%
\begin{tabular}{@{}lrrrrr@{}}
\toprule
 & tendance & vol. courte & lissé & teneur de marché & momentum\\
\midrule
rendement annuel & 14.4~\% & 11.6~\% & 7.2~\% & 22.8~\% & 1.1~\%\\
volatilité annuelle & 14.6~\% & 10.6~\% & 7.4~\% & 5.1~\% & 9.1~\%\\
ratio de Sharpe & 1.00 & 1.09 & 0.99 & 4.04 & 0.17\\
\quad erreur type, iid & 0.45 & 0.45 & 0.46 & 0.45 & 0.32\\
\quad erreur type, HAC & 0.46 & 0.45 & 0.63 & 0.52 & 0.31\\
\omterm{def:rs:performance-measurement:drawdown}{drawdown maximal} & $-33.6$~\% & $-14.8$~\% & $-11.0$~\% & $-4.2$~\% & $-34.0$~\%\\
drawdown le plus long & 791 j & 48 j & 22 m & 72 j & 2\,079 j\\
\omterm{def:rs:performance-measurement:drawdown}{ratio de Calmar} & 0.43 & 0.78 & 0.66 & 5.39 & 0.03\\
\omterm{def:rs:performance-measurement:sortino}{ratio de Sortino} & 1.48 & 1.87 & 2.00 & 6.31 & 0.24\\
\omterm{def:rs:performance-measurement:sortino}{ratio Oméga} en zéro (\omterm{def:rs:performance-measurement:hit}{profit factor}) & 1.17 & 1.49 & 2.07 & 2.07 & 1.03\\
asymétrie & $-0.04$ & 5.38 & 0.41 & $-0.66$ & 0.00\\
aplatissement & 3.0 & 102.5 & 2.8 & 9.9 & 3.4\\
\omterm{def:rs:performance-measurement:hit}{taux de réussite} & 53~\% & 76~\% & 58~\% & 65~\% & 51~\%\\
\omterm{def:rs:performance-measurement:beta}{bêta de marché} (erreur type HAC) & 0.00 (0.03) & 0.45 (0.09) & & 0.04 (0.01) & 0.00 (0.01)\\
\bottomrule
\end{tabular}
\end{center}

\textbf{Pour aller plus loin.} Doubler le notionnel de la volatilité courte et comparer le ratio de Sharpe affiché avec la probabilité d'un drawdown de 30~\%~; lisser le flux de tendance avec le profil $(0.5, 0.3, 0.2)$ et constater que son ratio de Sharpe augmente d'environ 62~\%.
\end{tutorial}



\section{Construction~: le module de performance}\label{bld:rs:performance-measurement:perf}

\begin{build}
\bfield{Rôle} Une fonction lit n'importe quel \texttt{BacktestResult} (ou flux de rendements) et fournit les indicateurs de ce chapitre avec leur incertitude, afin que
aucune stratégie ne soit jugée sur son seul ratio de Sharpe.
\bfield{Interface} \texttt{from\_returns(r, periods)}, \texttt{sharpe}, \texttt{lo\_sharpe(r, q, lags)}, \texttt{drawdown}, \texttt{max\_drawdown},
\texttt{drawdown\_spells}, \texttt{longest\_drawdown}, \texttt{calmar}, \texttt{sortino(r, periods, mar)}, \texttt{omega(r, threshold)}, \texttt{hit\_rate},
\texttt{profit\_factor}, \texttt{holding\_period(res)}, \texttt{alpha\_beta(r, bench, periods, lags)}, \texttt{smoothing\_profile(r, k)},
\texttt{unsmooth(r, theta)}, \texttt{tear\_sheet(res, bench, name, periods)}.
\bfield{Règles} Tout ratio de Sharpe est affiché avec son erreur type, en version iid et en version HAC~; les rendements mensuels ou lissés sont annualisés en tenant compte de leurs autocorrélations~;
les probabilités de drawdown proviennent d'un modèle de la stratégie, et non du maximum observé dans l'historique~; les statistiques de forme sont étiquetées comme celles de l'échantillon.
\bfield{Tests d'acceptation} \texttt{code/firm/perf/tests/}~: drawdowns, épisodes, Calmar, Sortino, Omega, \omterm{def:rs:performance-measurement:hit}{taux de réussite} et \omterm{def:rs:performance-measurement:hit}{profit factor} calculés à la main~; ratio de Lo égal à
$\sqrt q\,\mathrm{SR}_1$ sur des rendements indépendants, et correction d'un flux lissé~; profil de lissage retrouvé à 0.05 près et rendements vrais retrouvés
exactement après la période de démarrage~; un bêta et un alpha connus~; une durée de détention déduite d'un taux de rotation connu.
\bfield{Pour aller plus loin} Ratios de Sharpe déflatés à partir du \omterm{def:rs:the-research-process:log}{nombre d'essais} du \omterm{def:rs:the-research-process:log}{journal de recherche} (chapitre~20)~; fiches de performance glissantes~; la mesure insensible à la manipulation.
\end{build}

\begin{omsources}
\item A.~W.~Lo, ``The statistics of Sharpe ratios'', \emph{Financial Analysts Journal} 58(4), 2002.
\item M.~Getmansky, A.~W.~Lo and I.~Makarov, ``An econometric model of serial correlation and illiquidity in hedge fund returns'', \emph{Journal of Financial
Economics} 74(3), 2004.
\item M.~Magdon-Ismail, A.~F.~Atiya, A.~Pratap and Y.~S.~Abu-Mostafa, ``On the maximum drawdown of a Brownian motion'', \emph{Journal of Applied Probability}
41(1), 2004.
\item F.~Sortino and L.~Price, ``Performance measurement in a downside risk framework'', \emph{Journal of Investing} 3(3), 1994; PerformanceAnalytics (R),
documentation of \texttt{SortinoRatio}, \texttt{CalmarRatio} and \texttt{Omega}.
\item W.~F.~Sharpe, ``Mutual fund performance'', \emph{Journal of Business} 39(S1), 1966; ``The Sharpe ratio'', \emph{Journal of Portfolio Management}
21(1), 1994.
\item W.~Goetzmann, J.~Ingersoll, M.~Spiegel and I.~Welch, ``Portfolio performance manipulation and manipulation-proof performance measures'', \emph{Review of
Financial Studies} 20(5), 2007.
\end{omsources}

\section{Exercices}

\begin{exercise}[$\star$]\label{exo:rs:performance-measurement:1}
Une stratégie affiche un ratio de Sharpe de 1.0 sur cinq ans de rendements quotidiens. Calculer son erreur type et son intervalle à 95~\%.
\end{exercise}

\begin{exercise}[$\star$]\label{exo:rs:performance-measurement:2}
Le flux de tendance rapporte 14.4~\% par an avec un \omterm{def:rs:performance-measurement:drawdown}{drawdown maximal} de 33.6~\%. Calculer son \omterm{def:rs:performance-measurement:drawdown}{ratio de Calmar}. Pourquoi le \omterm{def:rs:performance-measurement:drawdown}{ratio de Calmar} d'un historique unique est-il instable~?
\end{exercise}

\begin{exercise}[$\star$]\label{exo:rs:performance-measurement:3}
Montrer que le \omterm{def:rs:performance-measurement:hit}{profit factor} est égal au \omterm{def:rs:performance-measurement:sortino}{ratio Oméga} au seuil zéro.
\end{exercise}

\begin{exercise}[$\star\star$]\label{exo:rs:performance-measurement:4}
Des rendements vrais de ratio de Sharpe 0.6 sont publiés lissés avec le profil $(0.5, 0.3, 0.2)$. Donner le facteur sur la variance, sur la volatilité, le
ratio de Sharpe publié et les deux premières autocorrélations.
\end{exercise}

\begin{exercise}[$\star\star$]\label{exo:rs:performance-measurement:5}
Avec les autocorrélations de l'exercice~4 (et nulles au-delà du retard 2), calculer le facteur d'\omterm{def:rs:performance-measurement:annualisation}{annualisation} de Lo pour des rendements mensuels et le comparer à $\sqrt{12}$. De combien
l'\omterm{def:rs:performance-measurement:annualisation}{annualisation} naïve surestime-t-elle le ratio de Sharpe~?
\end{exercise}

\begin{exercise}[$\star\star$]\label{exo:rs:performance-measurement:6}
Pourquoi un bêta de 0.45 sous-estime-t-il l'exposition de la stratégie de volatilité courte à un krach~? Utiliser les bêtas des jours de baisse et des jours de hausse.
\end{exercise}

\begin{exercise}[$\star\star\star$]\label{exo:rs:performance-measurement:7}
\emph{Programmation.} Avec \texttt{rs\_perf.dd\_probability}, donner la probabilité d'un drawdown de 30~\% sur un an, puis sur dix ans, pour chaque stratégie, et expliquer pourquoi
la probabilité de la stratégie de volatilité courte part de si haut.
\end{exercise}

\begin{exercise}[$\star\star\star$]\label{exo:rs:performance-measurement:8}
\emph{Trouver la faille.} «~Le \omterm{def:rs:performance-measurement:sortino}{ratio de Sortino} de notre fonds est de 1.9 et il n'a jamais eu de trimestre perdant en cinq ans~; son risque de baisse est donc inférieur à celui du
fonds de tendance voisin, dont le \omterm{def:rs:performance-measurement:sortino}{ratio de Sortino} est de 1.5 et qui a perdu un tiers de sa valeur.~»
\end{exercise}

\section{Problème~: deux ratios de Sharpe pour un seul}

\begin{problem}[{Problème du week-end --- deux ratios de Sharpe pour un seul}]\label{pb:rs:performance-measurement:1}
Les quatre flux du chapitre et le livre de momentum du chapitre~16, à travers \texttt{firm.perf}.

\medskip\noindent\textbf{Partie I --- Le ratio de Sharpe.}
\begin{enumerate}
\item Quels sont les ratios de Sharpe des historiques de tendance et de volatilité courte, et leurs erreurs types~?
\item Pourquoi quatre des cinq flux ont-ils presque la même erreur type~?
\item Que nous dit l'erreur type HAC sur le livre lissé~?
\item Quel est le ratio de Sharpe du teneur de marché et son erreur type, et que ne dit-il pas~?
\end{enumerate}

\medskip\noindent\textbf{Partie II --- Lissage.}
\begin{enumerate}[resume]
\item Quelles sont les volatilités publiée et vraie du livre lissé~?
\item Quels sont ses ratios de Sharpe publié, ajusté de Lo, non lissé et vrai~?
\item Quel profil de lissage est estimé, à partir de quelles autocorrélations~?
\item Pourquoi l'\omterm{def:rs:performance-measurement:annualisation}{annualisation} naïve est-elle fausse pour ce livre~?
\end{enumerate}

\medskip\noindent\textbf{Partie III --- Drawdowns et forme.}
\begin{enumerate}[resume]
\item Donner le \omterm{def:rs:performance-measurement:drawdown}{drawdown maximal} et la \omterm{def:rs:performance-measurement:drawdown}{durée du drawdown} la plus longue de chaque historique affiché.
\item Quelle stratégie a les meilleurs \omterm{def:rs:performance-measurement:drawdown}{ratios de Calmar}, Sortino et Omega dans les historiques affichés, et laquelle a la pire queue de distribution~?
\item Quelles sont l'asymétrie et l'aplatissement du flux de volatilité courte dans l'historique et sur 100 ans~?
\item Quelle part de ses historiques de cinq ans ne comporte aucun trimestre perdant~?
\item Quelles sont la durée de détention et le taux de rotation du livre de momentum~?
\end{enumerate}

\medskip\noindent\textbf{Partie IV --- Le verdict.}
\begin{enumerate}[resume]
\item Énoncer le \emph{résultat nommé}~: la probabilité d'un drawdown de 30~\% sur cinq ans pour chaque stratégie, et l'erreur type du ratio de Sharpe sous
l'autocorrélation du livre lissé.
\item Que donnent 200 années simulées pour le ratio de Sharpe et le \omterm{def:rs:performance-measurement:drawdown}{drawdown maximal} de chaque stratégie~?
\item Quels sont le bêta de la stratégie de volatilité courte, ses bêtas des jours de baisse et des jours de hausse~?
\item Auquel des deux allouer des capitaux, et comment en dimensionner l'allocation~?
\item Que devrait afficher chaque fiche de performance à côté d'un ratio de Sharpe~?
\item Comment un gérant pourrait-il contourner les mesures de ce chapitre~?
\item En une phrase~: que dit un ratio de Sharpe de un sur cinq ans~?
\end{enumerate}
\end{problem}

\section{Questions d'entretien}

\begin{interviewq}[$\star$ \iqroles{chercheur, trader}]\label{iq:rs:performance-measurement:1}
Une stratégie affiche un ratio de Sharpe de 1 sur trois ans de données quotidiennes. Quel degré de confiance accordez-vous à l'idée qu'il soit positif~?
\end{interviewq}

\begin{interviewq}[$\star\star$ \iqroles{chercheur, risque}]\label{iq:rs:performance-measurement:2}
Comment une stratégie peut-elle augmenter son ratio de Sharpe sans ajouter de compétence~? Comment le détecteriez-vous~?
\end{interviewq}

\begin{interviewq}[$\star\star$ \iqroles{chercheur}]\label{iq:rs:performance-measurement:3}
Les rendements mensuels d'un fonds ont une autocorrélation d'ordre 1 de 0.6. Qu'est-ce que cela change à son ratio de Sharpe, et comment le corrigez-vous~?
\end{interviewq}

\begin{interviewq}[$\star\star$ \iqroles{risque}]\label{iq:rs:performance-measurement:4}
Comment le \omterm{def:rs:performance-measurement:drawdown}{drawdown maximal} espéré d'une stratégie croît-il avec l'horizon~? Pourquoi le \omterm{def:rs:performance-measurement:drawdown}{drawdown maximal} d'un \omterm{def:rs:vectorised-backtests:backtest}{backtest} est-il une mauvaise estimation du risque~?
\end{interviewq}

\begin{interviewq}[$\star\star$ \iqroles{trader}]\label{iq:rs:performance-measurement:5}
Un teneur de marché gagne sur 65~\% des jours, avec un ratio de Sharpe de 4. Que voudriez-vous encore savoir~?
\end{interviewq}

\begin{interviewq}[$\star\star\star$ \iqroles{chercheur}]\label{iq:rs:performance-measurement:6}
Démontrer l'erreur type du ratio de Sharpe pour des rendements normaux indépendants.
\end{interviewq}
